% ============================================================================
% When Safe Search Budgets Do Not Transfer Across Graph-ANNS Rebuilds
% Anonymous ICLR submission — revised manuscript (v2, P0–P7 evidence integrated)
% ----------------------------------------------------------------------------
% Compile with the official ICLR kit: place iclr2027_conference.sty,
% iclr2027_conference.bst (and the ICLR banner files, e.g. iclr2027_logo.pdf)
% next to this file, then run:  pdflatex -> bibtex -> pdflatex x2
% If the ICLR style kit is temporarily unavailable, uncomment the FALLBACK
% block below to compile with plain article class.
% Figures required (same directory): fig1_setting.png, fig2_workflow.png,
% fig3_crossfamily.png, fig4_pooling.png, fig5_plane.png, fig6_ablation.png,
% fig7_ladder.png  (shipped in the artifact media/ folder).
% ============================================================================
\documentclass{article}

% ---- ICLR official style if present; otherwise self-contained fallback ----
\IfFileExists{iclr2027_conference.sty}{%
  \usepackage{iclr2027_conference,times}
}{%
  \usepackage[margin=1in]{geometry}
  \usepackage{times}
  \usepackage{natbib}
  \usepackage{fancyhdr}
  \pagestyle{fancy}\fancyhf{}
  \fancyhead[R]{\small\itshape Safe Search Budgets Across Graph-ANNS Rebuilds}
  \fancyfoot[C]{\small Under review as a conference paper at ICLR 2027 --- page \thepage}
  \renewcommand{\headrulewidth}{0pt}
}

\usepackage[utf8]{inputenc}
\usepackage[T1]{fontenc}
\usepackage{amsmath,amssymb,amsthm,mathtools}
\usepackage{booktabs}
\usepackage{graphicx}
\usepackage{multirow}
\usepackage{tabularx}
\usepackage{microtype}
\usepackage[hidelinks]{hyperref}

\newtheorem{theorem}{Theorem}
\newtheorem{proposition}[theorem]{Proposition}
\newtheorem{corollary}[theorem]{Corollary}

\newcommand{\DistComp}{\mathrm{DistComp}}
\newcommand{\TV}{\mathrm{TV}}
\newcommand{\KL}{\mathrm{KL}}
\renewcommand{\Pr}{\mathop{\mathrm{Pr}}}   % \Pr is a kernel command: renew, not new
\newcommand{\E}{\mathbb{E}}
\newcommand{\Acal}{\mathcal{A}}
\newcommand{\Ecal}{\mathcal{E}}
\newcommand{\BOT}{\bot}

\title{When Safe Search Budgets Do Not Transfer Across Graph-ANNS Rebuilds}

\author{Anonymous Authors\\
Paper under double-blind review\\
\texttt{anonymized artifact accompanies the submission}}

\begin{document}

\maketitle

\begin{abstract}
Graph approximate nearest-neighbor systems rebuild indexes after data refreshes, replica creation, and
randomized construction, yet downstream search-budget policies are commonly reused across builds. We
study when this transport is statistically justified, modeling each serialized build as an algorithmic
environment with implementation-native actions, an endpoint-aware failure event, and a categorical
unresolved state. A transcript-limited lower bound shows that build-blind policies must pay in one of
four currencies---unsafe execution, conservative computation, fallback, or abstention---and we
instantiate its premise: registered build pairs are near-chance distinguishable through hit-count
transcripts (TV $\le0.25$) while their per-query budgets conflict on 45--57\% of queries. A
fixed-target recovery theorem is audited per target: certification power, not margin existence, binds.
Across three dataset families, four scales (100K--10M), and three implementations, rebuilds cause
17.2--23.6\% incremental transport risk. We give a certified deployable policy (exchangeable-build
conformal pooling, Theorem~3): repeated queries served from cached labels carry a finite-sample risk
bound $\le\alpha$ at $1.32$--$1.44\times$ oracle cost, validated across every non-vacuous
$(\alpha,k)$; cold queries fall back to the maximum action. The deterministic rebuild contract
eliminates measured risk at $1.71$--$3.52\times$ build overhead. A learned per-query predictor does not
transfer ($R^2{=}0.015$). A cost-ordered decision rule over these
routes---contract, target re-profiling, conformal pooling, certify, abstain---is the paper's
actionable output.
\end{abstract}

\section{Introduction}
\label{sec:intro}

Graph approximate nearest-neighbor search (Graph-ANNS) is often tuned as though an index were a stable
substrate: after a rebuild with unchanged data and nominal hyperparameters, practitioners reuse budgets,
query-level predictors, and stopping rules calibrated on the previous index. The premise---that such
policies remain safe under a new realization of the construction algorithm---is rarely stated, and
reconfiguration violates it through input order, parallel scheduling, randomized state, or toolchain even
when the query law is unchanged. The reconfiguration arises in ordinary operation: serving replicas are
(re)created with different insertion orders and schedules, autoscaling adds replicas, failure recovery and
blue-green deploys rebuild from the same data, and upgrades rebuild under new toolchains; a policy
calibrated on one replica (a canary) is then reused across the fleet. We study the risk of this plausible
shortcut, make no claim that every operator reuses budgets, and condition all conclusions on registered
builds. The issue is not a mean-recall drop: builds with similar aggregate quality assign different
minimum safe actions to the same query, so reuse under-budgets the target while a uniformly larger action
pays conservative computation; and since primitive profiling is sub-second at 100K--1M, the gap is the
validity and decision semantics of transport, not profiling cost.

We formulate rebuild reconfiguration as a finite environment--action problem whose directed transport
separates unsafe execution, conservative computation, certified fallback, and abstention
(Section~\ref{sec:formulation})---an auditable risk--cost object. Our evidence addresses three layers: the
query-level safe-budget \emph{response} is strongly non-portable (17.17--23.60\% incremental risk at
100K, 22.16\% at 10M); the \emph{global policy} layer has a constructive answer (pooling 22 source builds
cuts risk below 1\%, certified under exchangeability by Theorem~3); and a \emph{learned per-query
predictor} does not transfer (Section~\ref{sec:scale}). Our theory separates impossibility from
conditional recovery: a transcript-limited lower bound whose premise we measure (hit-count TV $\le0.25$;
joint runtime features at most 0.72 separation; 45--57\% per-query conflict), and a fixed-target recovery
theorem whose five conditions we audit per target, locating the binding constraint at certification power
rather than margin existence. The testing inequalities are classical; the contribution is their
build-conditioned formulation, operational composition, and measured boundary.

We instantiate the theory in ICBA, a certified audit rather than a new traversal algorithm, and derive a
cost-ordered decision rule: enforce a verified deterministic contract when feasible; otherwise
re-profile the target with per-query truth labels when truth is affordable; otherwise deploy conformal
pooled replay when operational history permits (Theorem~3); otherwise collect target evidence and
certify a frozen action; otherwise abstain. Our contributions:
\begin{itemize}\itemsep1pt
\item A formalization of rebuild transport with implementation-native actions, categorical unresolved
endpoints, four deployment outcomes, and formal estimands (incremental risk, variation families).
\item Two complementary theories with measured boundaries---a transcript-limited lower bound with an
instantiated premise, and a fixed-target recovery theorem whose margin condition systematically
fails---plus a certified deployable policy: exchangeable-build conformal pooling (Theorem~3), which
upgrades pooling from bootstrap to finite-sample risk control.
\item Evidence across three dataset families (SIFT, Arxiv-Nomic, GloVe-100 angular), three geometric
spaces (L2-128, cosine-768, angular-100), four scales (10K farm to 10M), and three implementation
families (hnswlib, clean Faiss, Vamana), plus a mechanistic contract ablation, an economics ledger, and
a falsification ladder organizing negative method evidence.
\end{itemize}

\begin{table}[t]
\centering
\caption{Takeaway: deployment scenario $\rightarrow$ recommended route $\rightarrow$ measured boundary (all numbers registered-family conditional).}
\label{tab:takeaway}
\footnotesize
\setlength{\tabcolsep}{3pt}
\begin{tabular}{p{3.2cm}p{3.9cm}p{6.5cm}}
\toprule
Scenario & Route & Measured boundary \\
\midrule
Rebuild feasible offline; near-zero risk required & Deterministic contract (order + seed + single thread) & risk 0 at 100K/1M/10M (byte-identical); build $1.71$--$3.52\times$ (100K-registered; 10M: 2{,}945\,s); order pinning alone halves variation free \\
$\ge 1/\alpha-1$ prior builds with cached labels (19 at $\alpha{=}.05$, 9 at $\alpha{=}.10$) & Conformal pooled replay (Thm.~3 quantile) for repeat queries; max-action fallback for cold & certified $\le\alpha$ at every non-vacuous $(\alpha,k)$ across 3 scales $\times$ 2 implementations; DistComp $1.32$--$1.44\times$; uncached pooling degrades with scale (10M: 23.3\% at $k{=}7$) \\
Cost-sensitive, risk tolerance moderate & Canonical order pinning only & free ($0.80$--$0.93\times$ build); residual variation 30--34\% \\
Fresh target, budget for truth & Certify a frozen action (Thm.~2) or abstain & only max budget clears ($P{=}0.47$--$0.63$ at $m{=}59$); certification power, not margin, binds \\
No route yields a certified action & Abstain & safety preserved; regret bounded in Table~\ref{tab:constructive} \\
\bottomrule
\end{tabular}
\end{table}


\begin{figure}[t]
\centering
\includegraphics[width=.8\linewidth]{fig1_setting.png}
\caption{Build reconfiguration induces a policy-transport problem. A source-calibrated action may
under-budget the target build, while a uniformly large action may be safe but conservative. The build is
hidden only relative to the available probe transcript; ICBA separates measurement, certification, and
value.}
\label{fig:setting}
\end{figure}

Figure~\ref{fig:setting} summarizes the setting.

\section{Related Work}
\label{sec:related}

\paragraph{Positioning.}
Graph-ANNS construction and search: HNSW \citep{malkov2020hnsw}, NSG \citep{fu2019nsg},
Vamana/DiskANN \citep{subramanya2019diskann}, Faiss \citep{johnson2021faiss}, incremental maintenance
\citep{xu2022maintenance}, construction cost and pruning theory \citep{prokhorenkova2020graph}, order
effects on recall \citep{elliott2024order}, per-index budget autotuning \citep{bae2026qbat}.
Query-adaptive budgets and cost estimators operate per profiled index
\citep{li2020learned,chatzakis2025darth,horchidan2025conann,wang2026annie,zhang2026exploration}. Risk
control and safe selection: LTT, conformal risk control, selective classification, safe BAI
\citep{angelopoulos2025ltt,bates2021rcps,angelopoulos2024crc,elyaniv2010selective,wang2022safebai,garivier2016bai}.
Distribution shift treats environments as sampling units
\citep{bendavid2010theory,johansson2019support,duchi2025multienv}. None of these lines asks whether a
safety-critical budget transports across an independent rebuild of the same data; that transport-validity
question, its measured boundaries, and a certified deployment policy are this paper's subject (full
discussion in Appendix~\ref{app:related}).

\section{Theory}
\label{sec:theory}

\subsection{Main Theorem I: transcript-limited non-portability}

Consider two build environments $E_0,E_1$ where a procedure observes $T\sim P_i^{T}$ and outputs an
action through any randomized Markov kernel measurable in $T$. Let $D_i$ be the region of acceptable
actions in $E_i$, assume $D_0\cap D_1=\varnothing$, and let $l_i\ge\Delta\,\mathbf{1}\{a\notin D_i\}$.

\begin{theorem}[transcript-limited lower bound]
Every transcript-based randomized procedure satisfies
\begin{equation}
\max_{i\in\{0,1\}}\E[l_{i}]\ \ge\ \frac{\Delta}{2}\bigl(1-\TV(P_{0}^{T},P_{1}^{T})\bigr).
\label{eq:thm1}
\end{equation}
\end{theorem}

If the transcript under the registered probe budget cannot identify which region applies, no
transcript-measurable policy can uniformly avoid unsafe execution, conservative computation, fallback, or
abstention. For an adaptive transcript with KL divergence $K_m$, Bretagnolle--Huber and Pinsker give
\begin{equation}
\max_{i}\E[l_{i}]\ \ge\ \frac{\Delta}{4}e^{-K_{m}},\qquad
\max_{i}\E[l_{i}]\ \ge\ \frac{\Delta}{2}\max\!\Bigl(0,\,1-\sqrt{K_{m}/2}\Bigr).
\label{eq:bh}
\end{equation}
The tools are classical \citep{yu1997assouad}; the contribution is the reduction from a wrong build
decision to the four Graph-ANNS deployment outcomes.

\paragraph{Probe-class evidence on the premise (Section~\ref{sec:scale}).}
The premise concerns the policy's \emph{full} observable $(q,T)$: marginal transcript indistinguishability
alone would not imply joint indistinguishability, and a classifier's accuracy bounds TV from
\emph{below} ($\TV\ge 2A{-}1$), which cannot by itself yield the loss bound. The valid object is a
\emph{confidence upper bound} on TV for a declared transcript class. For hit-count transcripts, a DKW
concentration argument on the finite transcript support gives $\TV\le0.39$ at 95\% confidence
(plug-in 0.25 $+\,0.14$ concentration), so Theorem~1 yields a valid loss lower bound of
$0.305\,\Delta$ for that class; joint runtime features separate more (worst pair accuracy 0.72), which
\emph{weakens} the bound for that class ($\TV\ge0.44$) and is reported as a risk-inflation diagnostic,
not as premise evidence. The theorem remains conditional on the declared probe class; the 45--57\% conflict rate and the
common-safe-action analysis (21.9\% no-conflict-free transport, Section~\ref{sec:constructive}) supply
the remaining instantiation.

\subsection{Main Theorem II: finite-class conditional recoverability}

Fix a named target $E_t$, a frozen family $\Acal=\{a_1,\dots,a_M\}$, a declared $P_Q$, and the
endpoint-aware loss, with: (i) an action of risk $\le\delta-2\gamma$; (ii) action-relevant
identifiability; (iii) positive margin; (iv) independent selection/certification roles; and (v) a valid
fallback or abstention. On the simultaneous estimation event
\begin{equation}
\mathcal{G}
=\Bigl\{\sup_{a\in\Acal}|\hat r_{a}-r_{a}|\le\varepsilon_{R}\Bigr\}
\cap\Bigl\{\sup_{a\in\Acal}|\hat C_{a}-C_{a}|\le\varepsilon_{C}\Bigr\},
\label{eq:eventG}
\end{equation}
screening by $\hat r_{a}+\varepsilon_{R}\le\delta$ and certifying the frozen choice on independent data
gives
\begin{equation}
\Pr\bigl[r_{E_{t}}^{\hat a}\le\delta\ \ \text{or a valid fallback or abstention}\bigr]\ \ge\ 1-\alpha.
\label{eq:thm2}
\end{equation}
This is a sufficient protocol plus necessary obstructions, not an if-and-only-if characterization.
\begin{theorem}[fixed-target conditional recoverability]
\label{thm:recover}
Under conditions (i)--(v),
\begin{equation}
\Pr\bigl[r_{E_{t}}^{\hat a}\le\delta\ \ \text{or a valid fallback or abstention is invoked}\bigr]\ \ge\ 1-\alpha.
\label{eq:thm2env}
\end{equation}
\end{theorem}
Section~\ref{sec:constructive} audits (i)--(v) per target: (iii) passes point-estimate-wise at small
margins, but margin-band actions carry 3--4\% risk and fail zero-failure certification at practical sample
sizes---certification power binds, not margin existence.

\subsection{Exchangeable-build conformal pooling (Theorem 3)}
\label{sec:thm3}

Theorem~2 certifies a \emph{fixed} action and needs a margin; pooling certifies a \emph{per-query order
statistic} and needs only that the target and source builds are exchangeable units (the query law is
arbitrary; queries need not be independent).

\begin{theorem}[exchangeable-build conformal pooling]
Let $B_{E_1}(q),\dots,B_{E_k}(q),B_{E_t}(q)$ be exchangeable across builds for each fixed query $q$.
Let $m=\lceil(1-\alpha)(k+1)\rceil$ and deploy, per query, the $m$-th smallest of the source minimum safe
actions, where the unresolved symbol $\bot$ participates in the order as $+\infty$ (a pool whose $m$-th
smallest is $\bot$ yields abstention). If $m\le k$ (equivalently $\alpha\ge 1/(k+1)$), then
\begin{equation}
\Pr\bigl[B_{E_t}(q)>a(q)\bigr]\ \le\ \alpha.
\label{eq:thm3}
\end{equation}
\end{theorem}

\begin{proof}
Under exchangeability the rank of $B_{E_t}(q)$ among the $k{+}1$ values is uniform on
$\{1,\dots,k{+}1\}$; the deployment event is $\{\mathrm{rank}>m\}$, of probability at most
$(k{+}1-m)/(k{+}1)\le\alpha$.
\end{proof}

The guarantee bounds \emph{label coverage} $P[B_{E_t}(q)>a(q)]\le\alpha$; deployment safety further
requires that deploying any $a\ge B_{E_t}(q)$ succeeds. This bridge is \emph{measured, not assumed}: the
non-monotone violation rate, conditioned on being at or above the target's own minimum safe point
($\nu=\Pr[Z_{E_t}(q,a){=}1 \mid a \ge B_{E_t}(q)]$ over all (query, build, $a \ge B_{E_t}(q)$)
triples, exhaustively checked), is $\nu=0$ on hnswlib-100K (both datasets), Vamana (both stages), and
10M, and $\nu=0.39\%$ on 1M. Because conformal-pooling deployments satisfying $a(q) \ge B_{E_t}(q)$
lie in this condition, execution failure is bounded by $\alpha+\nu$ ($\le 5\%$ at the certified
operating point, measured 0.96--1.35\%); the residual risk arises solely from transported actions
\emph{below} the target's minimum safe point ($B_s(q) < B_{E_t}(q)$)---the transport phenomenon
itself, not a monotonicity violation. A stricter tail-label variant ($B^{\mathrm{tail}}$ = min action after which all
registered actions succeed) gives execution safety by construction (realized 0.13--0.18\%) at halved
cache coverage---both variants are reported. Measured boundaries: the certificate resolves only
$\alpha\ge 1/(k+1)$ (a 2\%-class certificate needs $k\ge 49$), and target-correlated source selection
inflates realized risk $3.3\times$ (4.43\% vs.\ 1.31\% at $\alpha{=}0.10$, $k{=}9$), within the
nominal envelope on this sample---an exchangeability-violation diagnostic (the premise, not the certificate bound, is violated).
The \emph{tiered} deployment composes the certificate with the measured fallback: cached actions for
repeated queries, maximum action for cold queries, giving
$r_{\mathrm{tiered}}\le\rho\,\alpha+(1-\rho)\,r_{\max}$ at repeat rate $\rho$; meeting $\delta$
requires $\rho\ge(r_{\max}-\delta)/(r_{\max}-\alpha)$, satisfied at every measured $\rho$ at 100K
($r_{\max}{=}1.3\%$; interpolation 1.26\%$\to$0.09\%, Section~\ref{sec:constructive}).

\begin{table}[t]
\centering
\caption{Theorem-3 validity across scales and implementations (leave-one-build-out; 100K rows use 24
targets $\times$ 12 draws, farm rows 30 draws). Valid at every non-vacuous $(\alpha,k)$; the
weighted-selection row shows the exchangeability premise is load-bearing. DistComp = distance
computations vs.\ per-query oracle.}
\label{tab:thm3}
\footnotesize
\setlength{\tabcolsep}{3.4pt}
\begin{tabular}{llcccc}
\toprule
Cell & $\alpha$ & $k$ & Realized failure & Valid & DistComp \\
\midrule
hnswlib SIFT & 0.05 & 23 & 0.96\% & \checkmark & $1.44\times$ \\
hnswlib Arxiv & 0.05 & 23 & 1.35\% & \checkmark & $1.32\times$ \\
Faiss SIFT (clean) & 0.05 & 23 & 0.05\% & \checkmark & --- \\
Faiss Arxiv (clean) & 0.05 & 23 & 0.05\% & \checkmark & --- \\
hnswlib SIFT & 0.10 & 9 & 1.31\% & \checkmark & $1.43\times$ \\
Faiss SIFT (clean) & 0.10 & 10 & 1.40\% & \checkmark & --- \\
SIFT-1M (hnswlib) & 0.125 & 7 & 0.92\% & \checkmark & --- \\
deep-10M (hnswlib) & 0.125 & 7 & 0.64\% & \checkmark & --- \\
10K farm & 0.05 & 49 & 0.72\% & \checkmark & --- \\
10K farm & 0.10 & 9 & 2.48\% & \checkmark & --- \\
hnswlib SIFT (corr.-weighted) & 0.10 & 9 & 4.43\% & $\times$ & $1.36\times$ \\
\midrule
\multicolumn{6}{l}{\emph{Vacuous configurations behave as the theorem predicts:}} \\
hnswlib SIFT & 0.05 & 5 & 3.50\% & (m$>$k) & --- \\
\bottomrule
\end{tabular}
\end{table}

\subsection{Certification scale and economic feasibility}

Uniform concentration gives the sufficient order $m=O(\log(M/\alpha)/\gamma^{2})$; zero-failure
Clopper--Pearson certification requires $m_{\min}=\lceil\log\alpha_{l}/\log(1-\delta)\rceil$---59 queries
for a single frozen action, 94 under Bonferroni over $M{=}6$, and pass probability $(1-r)^m$ at true risk
$r$ (the maximum-budget action passes only 47--63\% of the time at $m{=}59$). For economics,
$D=G_{\mathrm{online}}-\Pr(R)\,\Delta C_{f-c}$ and $N^{\star}=K_0/D$ if $D>0$: an accounting identity, not
an efficiency theorem.
\section{ICBA: Certified Audit and Actionable Mitigation}
\label{sec:icba}

ICBA (the artifact's name for this certified budget-audit workflow) receives registered source and target
builds, implementation-native actions, disjoint
design/selection/certification/evaluation roles, exact truth for authorized labeled roles, one
endpoint-aware failure event, and a frozen candidate family; it returns certified reuse, target-specific
recalibration, an independently justified fallback, or abstention---never treating the maximum registered
action as safe without evidence. It validates roles, replay, and action semantics; builds the categorical
safe-budget response without endpoint imputation; measures directed transport; freezes any
target-data-selected candidate before familywise-controlled certification; and audits mean, tail,
truth/control, fallback, and break-even quantities separately, leaving missing costs not estimable. A
complementary mitigation removes the build environment: the deterministic-rebuild contract fixes input
order, seed, threads, implementation, and toolchain with byte-level replay (Section~\ref{sec:contract}).
The resulting cost-ordered decision rule: (1) deterministic contract when feasible; (2) conformal pooled
replay when history permits (Theorem~3); (3) target certification, expecting only the maximum budget to
clear on grids like ours; (4) abstain.

\section{Experimental Protocol}
\label{sec:protocol}

We evaluate SIFT-100K and Arxiv-Nomic-100K at Recall@10 with hit requirement $h{=}10$ (all ten
neighbors returned). hnswlib and clean Faiss-100K cells have 24 registered builds per dataset and all 552
directed pairs; the Vamana stage has 12 builds and 36 preregistered directions; each 100K cell uses 750
evaluation queries; grids are native (hnswlib 10--200, Faiss 16--512) and never equated across systems.
The preregistered 1M and 10M replications mirror this design (8 builds; 500 queries; forensics gate on
query/base content overlap passed before any risk computation; contract identity re-verified at 10M).
Evaluation queries are content-hash disjoint from the indexed base and historical roles (all overlap
counts zero; Vamana caveats recorded in Appendix~\ref{app:crossfamily}). Primary intervals are 95\%
query-bootstrap (5{,}000 resamples, seed 991) over the full directed-pair vector, conditional on the
registered build family. Full protocol details are in Appendix~\ref{app:protocol}.

\section{Results}
\label{sec:results}

\subsection{Strictly comparable HNSW transport evidence}

The clean hnswlib and Faiss experiments agree on the main conclusion. Across the four data--implementation
cells, 75.87--91.07\% of evaluation queries exhibit finite-action safe-budget variation across registered
builds, whereas endpoint-state variation is only 0.53--2.67\%. Incremental source-to-target transport risk
ranges from 17.17\% to 23.60\% (Table~\ref{tab:crossfamily}). Every 95\% query-bootstrap interval excludes
the preregistered 2\% materiality threshold. Thus the primary signal is switching among finite native
actions, not a numerical artifact created by unresolved endpoints.

\begin{table}[t]
\centering
\caption{Registered cross-family evidence at Recall@10 with $h{=}10$. HNSW rows share the repaired common
estimand; the estimand crosswalk in Appendix~\ref{app:crossfamily} tabulates the coexisting registered
values.}
\label{tab:crossfamily}
\footnotesize
\setlength{\tabcolsep}{2.6pt}
\begin{tabular}{llccccc}
\toprule
Impl. & Dataset & B/p & $V_{\mathrm{fin}}$ & Risk [95\% CI] & $V_{\mathrm{end}}$ & Within-impl.\ cost \\
\midrule
hnswlib & SIFT-100K & 24/552 & 89.33\% & 21.55\% {\tiny[20.76, 22.33]} & 2.67\% & DistComp $+19.45\%$ \\
hnswlib & Arxiv-100K & 24/552 & 75.87\% & 17.17\% {\tiny[16.33, 18.06]} & 2.00\% & DistComp $+14.88\%$ \\
Faiss & SIFT-100K & 24/552 & 91.07\% & 23.60\% {\tiny[22.78, 24.39]} & 0.53\% & N/E (batch-cum.) \\
Faiss & Arxiv-100K & 24/552 & 75.87\% & 17.77\% {\tiny[16.86, 18.68]} & 1.87\% & N/E (batch-cum.) \\
\bottomrule
\end{tabular}
\end{table}

Robustness checks do not alter this conclusion. Removing the eight queries with the largest hnswlib risk
contribution leaves increments of 21.40\% and 16.90\%; deleting the top 1\% of Faiss risk contributors
leaves 23.43\% and 17.53\%; all 24 leave-one-build estimates remain positive on every cell. These are
robustness diagnostics conditional on the registered build family, not a population model for future
rebuilds.

\paragraph{Quality-threshold sensitivity.}
The $h{=}10$ event (all ten neighbors returned) is the strictest registered target, so we recomputed the
\emph{incremental} family transport risk (Eq.~\ref{eq:delta} caliber: transported minus target-bottom
reference, averaged over all directed pairs) under relaxed hit requirements $h\in\{8,9\}$ for every cell
(Appendix~\ref{app:protocol}, Table~\ref{tab:hsens}). By construction the $h{=}10$ column reproduces the registered incremental
estimates (the hnswlib rows land on the E4-primary caliber 21.57/17.12, which the Appendix~F crosswalk
tabulates alongside the repaired-common 21.55/17.17; Faiss and 1M match Table~\ref{tab:crossfamily} to
the digit). Risk decreases as the requirement relaxes, but the minimum over all cells at $h{=}8$ is
5.96\%---still three times the 2\% materiality gate. The phenomenon is not an artifact of near-exact
retrieval.

\subsection{Vamana-style supporting boundary}

The boundary panel (Table~\ref{tab:vamana}, appendix) reproduces the risk direction on both datasets under
its original stage-specific estimand (12 builds, 36 asymmetric directions, beam width one), broadening the
construction-family scope without supporting numerical equality to HNSW. Cost is equally non-universal:
hnswlib shows DistComp penalties of 19.45\%/14.88\%, the Vamana-style estimates sit near 1\% with
intervals crossing zero, and clean Faiss distance-computation cost is not estimable (batch-cumulative
counter).

\subsection{Deterministic rebuilding removes the registered environment}
\label{sec:contract}

We fix the input order, seed, thread count, implementation, and toolchain, and require byte-identical
serialization and search replay. Three repeated builds per dataset satisfy all identity checks and have
measured transport risk zero. Relative to the original multithreaded construction, the contract preserves
Recall@10 within 0.0004, leaves p95/p99 safe budgets unchanged, changes mean safe budgets by ratios 0.976
(SIFT) and 0.999 (Arxiv), and increases build time by $3.52\times$ and $1.71\times$. This is an actionable
environment-elimination contract---a hnswlib deterministic construction control---not a recovery algorithm
for changed data or software.

\paragraph{Chain ablation (descriptive).}
A four-tier chain---D0C (random order, 8 threads, random seeds), D1 (canonical order pinned), D2 (seed
pinned), D3 (single-threaded)---attributes the contract's effect: pinning the insertion order to a
canonical order halves minimum-safe-action variation (68.3 to 34.1\% on SIFT; 63.2 to 30.8\% on Arxiv) at
build-time ratios $0.80\times$ and $0.93\times$, i.e., free. Pinning the seed under 8-thread construction is
nearly inert ($-1.3$ and $-4.3$ points), and the six D2 builds with identical order and seed produce six
distinct serialized indexes: parallel scheduling is a residual nondeterminism source that the exposed seed
does not control. Single-threaded construction is therefore the necessary and sufficient registered clause
for byte identity (3/3 identical, variation 0) and carries the entire overhead ($4.73\times$ and
$1.86\times$ relative to the pinned 8-thread tier; $3.52\times$ and $1.71\times$ relative to D0C). The
ablation is a chain, not a full factorial; effects are descriptive. At 1M the contract remains feasible:
single-threaded builds took a median 252\,s (13 builds, range 237--267\,s; byte-identity rebuilds 238\,s),
a $19.5\times$ single-thread scaling from 100K for $10\times$ the data; the 8-thread 1M comparison was
not measured, so the $1.71$--$3.52\times$ ratios remain 100K-registered values.

\subsection{Constructive baselines and certification}
\label{sec:constructive}

We evaluate seven policies on the hnswlib cells by retrospective replay over the frozen per-query records
(roles 375/94/281 disjoint, seed 991; bootstrap as registered). Table~\ref{tab:constructive} reports
target-truth usage, risk under the four-outcome semantics, DistComp versus the per-query oracle, and
certification outcomes.

\begin{table}[t]
\centering
\caption{Constructive decision table (hnswlib, registered grid, retrospective replay). M2 dominates on risk
at $1.34$--$1.45\times$ oracle cost but requires a per-query source-label cache (see text) and is
uncertified; certified routes collapse to the maximum budget.}
\label{tab:constructive}
\scriptsize
\setlength{\tabcolsep}{2.8pt}
\begin{tabular}{p{3.2cm}p{1.0cm}p{3.0cm}p{1.4cm}p{4.0cm}}
\toprule
Method & Truth & Risk (SIFT / Arxiv) & DistC. & Certification \\
\midrule
M1 naive transport ($k{=}1$) & 0 & 21.89\% / 18.06\% & 2.6/2.7$\times$ & none \\
M2 pooled replay (max over 22 cached source labels) & 0 target$^{\dagger}$ & 0.18\% / 0.09\% (abst.\ 2.4/2.7\%) & 1.45/1.34$\times$ & uncert.; LOBO $\le$2.1\% \\
M3 frozen max-action CP ($M{=}1$, $m{=}59$) & 59 q & 0.87\% / 1.36\% & 4.16/5.28$\times$ & pass-prob mean 0.63/0.47; certifiable targets 71\%/38\% \\
M4a point-screen sel.-then-cert. & 469 q & --- & --- & 0/48 (zero-fail.\ violated) \\
M4b Theorem-2 screen, then certify & 469 q & --- & --- & 0/48; margin too thin \\
M5 always maximum action & 0 & 0.80\% / 1.26\% & 4.09/5.14$\times$ & not auto.\ safe \\
M2 pooled replay (max over 22 cached source labels) & 0.12--0.99\,s per 59 q (truth incl.) & endpoint mass only & 1.00$\times$ & uncert.; dominates M2 when truth affordable \\
M2 pooled replay (max over 22 cached source labels) & 0.12--0.99\,s per 59 q (truth incl.) & endpoint mass only & 1.00$\times$ & uncert.; dominates M2 when truth affordable \\
M5$'$ target re-profile + per-query truth cache & 0.12--0.99\\,s per 59 q (truth incl.) & endpoint mass only & 1.00$\\times$ & uncert.; dominates M2 when truth affordable \\
M6 per-query oracle & full & endpoint mass only & 1.00$\\times$ & non-deployable ref. \\\\
\bottomrule
\end{tabular}
\end{table}

\paragraph{Theorem-2 verdict and D-region instantiation.}
The disjoint-region premise holds for a substantial conflict subset. Over the frozen hit tables, 21.9\%
of directed transports (SIFT-100K) admit \emph{no} common safe action within the source action's own
cost (still 6.9\% allowing twice the target budget), with median action gap 30 (SIFT) / 20 (Arxiv),
$p_{90}$ 80/70---computed by scanning both hit tables upward from $\max(B_s,B_t)$ to find the first
action safe in both (method in the artifact).
Conditions (i), (ii), (iv), (v) are satisfiable on the registered grid; (iii) fails for all 48 targets:
every non-maximal action's risk lies inside the margin band, so Theorem-2 screening passes (it selects
$\mathrm{ef}{=}120$ on eight targets and $\mathrm{ef}{=}200$ on sixteen per dataset) but zero-failure
certification then fails on every target. The certified route collapses to
maximum-budget-or-abstain, and even the maximum action certifies only with probability 0.47--0.63 (per-target pass-probability
mean; the share of targets whose pass-probability exceeds 0.5 is 71\%/38\%,
Table~\ref{tab:constructive}) at $m=59$ because its residual risk (0.008--0.013, concentrated in
always-infeasible queries) violates the zero-failure requirement in roughly half the draws. The CP power curve locates this precisely: certifying
actions of risk 1--2\% needs 59--99 queries, but the margin-band actions carry 3--4\% risk and require
289--1{,}149 queries---the zero-failure protocol at registered sample sizes, not candidate nonexistence,
is what fails here. A $\gamma$-sensitivity audit sharpens the verdict: at
$\gamma{=}0.005$ a non-maximal action with selection-block risk $\le\delta-2\gamma$ \emph{exists} on
16/24 (SIFT) and 21/24 (Arxiv) targets (16/24 and 8/24 at $\gamma{=}0.01$; none at $\gamma{=}0.02$), so
the margin condition is not absolutely absent---what fails everywhere is the \emph{certification step}:
margin-band actions carry 3--4\% risk and pass zero-failure certification at $m{=}94$ with probability
$(1-r)^{94}\approx2$--$6\%$. Certification power, not margin existence, is the binding constraint at
registered sample sizes. This is a measured boundary of the recovery theorem on this grid---conditional,
not an impossibility claim. Pooling is the constructive winner when the serving workload repeats queries whose source-side safe
budgets are cached: M2 deploys the maximum over stored per-query source labels and needs no target truth,
but it is a \emph{replay policy with a label cache}---for cold queries it must fall back (e.g., to the
maximum action or abstention), and its uncertainty remains a registered-family bootstrap, not an
unseen-build guarantee. Deployability therefore equals the workload's query-repeat coverage, which is a
deployment property we measure rather than assume. Theorem~3 upgrades this picture from bootstrap to
certificate: under build exchangeability the $(1{-}\alpha)$-quantile pooling action bounds per-query
transport risk by $\alpha$ finitely---and the measured coverage (Table~\ref{tab:thm3}) confirms it
out-of-sample across a complete validity matrix: two implementations at 100K (hnswlib 0.96--1.35\%,
clean Faiss 0.05\% at $\alpha{=}0.05$, $k{=}23$; DistComp $1.32$--$1.44\times$), the 1M and 10M cells
at their pool-resolution limit ($k{=}7$, $\alpha{\ge}0.125$: realized 0.92\% and 0.64\%), the
200-build farm at certified depth ($k{=}49$, $\alpha{=}0.05$: 0.72\%), the Vamana family
($k{=}5$, $\alpha{\ge}1/6$: realized 0\%), and GloVe-100 at its resolution limit (6.47\%). The tiered composition (cache for repeats, maximum action for cold
queries) then meets $\delta$ at every measured repeat rate $\rho\in[0,1]$: effective risk interpolates
$1.26\%\to0.09\%$ (Arxiv) and $0.80\%\to0.18\%$ (SIFT) as $\rho$ goes $0\to1$ (Table~\ref{tab:constructive},
coverage column). An \emph{online} pooled policy that searches all sources for cold queries is not
serviceable: serving cost $1.73$\,s per cold query (23 graded source searches) versus $276\,\mu$s for the
maximum action---measured on the fresh-query workload. Non-max aggregation does not help: at $k{=}22$ the $0.9$-quantile of source
actions is nearly identical to the max (1.07\% vs.\ 0.96\% risk at 45.5\% vs.\ 48.4\% conservatism on
SIFT), while median pooling cuts conservatism to 16--22\% but leaves 13--21\% risk---the response
diameter is driven by outlier sources, so there is no cheaper order-statistic operating point.

\subsection{Scale, learned-predictor transfer, and distinguishability}
\label{sec:scale}

\paragraph{10M and GloVe-100 replications (preregistered).}
On ann-benchmarks deep-image (9.99M $\times$ 96, angular), eight builds in the registered
configuration and 500 preregistered queries (forensics gate passed before any risk
computation) give incremental transport risk \textbf{22.16\% [21.66, 22.68]} with
$V_{\mathrm{fin}}{=}87.2\%$ versus $V_{\mathrm{end}}{=}28.2\%$: the phenomenon is
scale-stable across three orders of magnitude (21.55--23.60\% at 100K, 21.87\% at 1M).
The deterministic contract is byte-identical at 10M (three same-configuration builds,
sha-identical across days) at a median 2{,}945\,s single-thread build; profiling remains
sub-second (0.091\,s for 59 queries $\times$ 6 actions). The pooling boundary amplifies
with scale: the uncached $k$-curve is nearly flat ($k{=}3$: 26.8\%, $k{=}5$: 24.5\%, $k{=}7$:
23.3\%), i.e., the source requirement has outgrown practical history at 10M, while
$h$-sensitivity stays flat (22.16--22.74\% for $h{=}10{..}8$). Theorem~3 remains valid at the 10M
pool resolution (realized 0.64\% $\le\alpha{=}0.125$ at $k{=}7$) with 36\% abstention from the
larger endpoint mass---the certificate survives, the cheap uncached pooling does not.

\paragraph{GloVe-100 (third dataset family, angular geometry).}
A preregistered eight-build cell on GloVe-100 (1.19M $\times$ 100, cosine/angular) gives incremental
transport risk \textbf{17.54\% [16.89, 18.11]} with $V_{\mathrm{fin}}{=}67.4\%$ versus
$V_{\mathrm{end}}{=}27.8\%$ and 28.4\% unresolved mass---the angular geometry is harder (endpoint
mass triples versus the L2 cells) yet finite-action switching remains the dominant channel,
$h$-sensitivity stays flat (17.54--19.66\% for $h{=}10{..}8$), and Theorem~3 holds at the pool
resolution (6.47\% $\le\alpha{=}0.125$ at $k{=}7$). Uncached pooling is at its worst here
($k{=}7$: 42.5\%): the third dataset confirms both the phenomenon's generality and the
scale/geometry-dependence of every mitigation.

A 200-build \emph{population farm} (two arms $\times$ 100 rebuilds on SIFT-10K: registered-style
single-thread and deployment-style 8-thread) measures the premises directly at population scale. Every one
of the 9{,}900 directed pairs exceeds the 2\% gate (mean 23.4\%, p95 26.6\%); the two arms are
statistically identical (23.42\% vs 23.45\%)---parallel-scheduling noise does not amplify population risk
beyond registered-style variation. Exchangeability itself is testable on the farm: the rank of a held-out
build's minimum safe action within random pools is uniform after randomized tie-breaking (KS median
0.114/0.118), and Theorem-3 coverage at $k{=}9$, $\alpha{=}0.10$ realizes 2.53\%/2.62\% $\le\alpha$.
The theorem's premise and guarantee both hold on an independent farm; a no-truth budget rule (queue-
saturation ef selection) degenerates to the maximum action (0.8--1.4\% risk at $3.8$--$5.1\times$ oracle),
closing the ``profile without truth'' loophole by measurement.

The predictor probe deserves emphasis: a regressor trained on 23 source builds predicts the target's
minimum safe action no better than chance ($R^2=0.015$), and even in-target cross-validation leaves
22--24\% risk---the deployment-time transcript of the cheapest action carries almost no information about
which build-specific neighborhood structure makes a larger budget necessary. The three-layer boundary is
now complete: response labels are strongly non-portable (layer 1), global policies have a constructive
bounded answer (layer 2, Section~\ref{sec:constructive}), and learned per-query predictors do not transfer
on these features (layer 3). We did not test richer probe features or model classes; the claim is
feature-class-conditional.

\subsection{Economics ledger and grid sensitivity}

Transport risk is grid-indexed: on registered subgrids (drop-min, drop-max, coarse 3-action) mean family
risk stays between 10.3\% and 23.8\%---always above the 2\% materiality gate---while coarser grids
mechanically lower measured risk because the minimum safe action can only grow. We therefore claim
phenomenon-level grid-robustness, not resolution-invariant risk levels, and keep native grids separate
across implementations.

\paragraph{End-to-end: transport versus profile-the-target.}
Because primitive profiling is sub-second, the decision rule's real competitions are ``rebuild, then
profile the target'' and ``profile the target and certify.'' Brute-force truth for a 500-query cache
costs $\sim$0.05\,s at 100K, $\sim$0.5\,s at 1M and $\sim$5\,s at 10M---cheap at every scale we
measure; what is \emph{not} affordable is per-query truth for a full production stream, which is where
the label cache (M2/M5$'$) and the certificate (Thm.~3) operate. Table~\ref{tab:e2e} assembles the full offline-cost / risk / certification
ledger from measured components: profiling-and-certifying is indeed cheap (0.12--0.99\,s per target,
truth included), but on this grid it can only \emph{certify the maximum budget} and passes zero-failure
certification on 47--63\% of draws---cost was never the binding constraint; candidate attainability and
certification power are. The deterministic contract remains the only zero-risk route, and the pooled
replay policy the only low-risk uncertified route, so the decision rule's ordering survives the cheap
profiling comparison.

\section{Theory-Guided Falsification Ladder}
\label{sec:ladder}

Sixteen attempted recovery routes are organized as a falsification ladder, each pruned at its first
failed prerequisite---action semantics, candidate attainability, observable information, structural
bridges, certification, tail behavior, or deployment value. Compressed verdicts: reuse/recalibration
retains fixed-target safety but no stable history increment (the predictor probe quantifies the negative
side); build selection finds no interval-supported jointly feasible action among 72 frozen
build-budget actions; portal rescue shows real complementarity (1{,}064/727 hits; 426/348 rescues) but no
truth-free trigger and a 79.5--93.4\% cost-compression requirement; the unified auditor is safe without
value (36/36 accepted decisions safe, all at maximum budget; regret 470.34/684.49). Family table, figure,
casebook, and provenance notes: Appendix~\ref{app:ladder}.

\section{Limitations and Scope}
\label{sec:limitations}

Primary evidence covers two 100K datasets and four HNSW cells plus single-implementation 1M/10M
replications (8 builds each); Vamana uses a different estimand and beam width one. We do not model future
builds, data refreshes, or open-world safety; Faiss distance-computation cost and Vamana wall-clock/I/O
are not estimable. The lower bound is probe-class-limited (hit-count and joint runtime features tested;
the conflict rate supports but does not instantiate the disjoint regions and gap $\Delta$), and the
recovery theorem binds at certification power on this grid; richer candidate families might change that.
The contract's ablation is a descriptive chain. The constructive policy is a retrospective replay:
pooling needs $k\ge 1/\alpha-1$ cached builds (19 at $\alpha{=}0.05$; 10M-scale uncached pooling
fails), its risk is an uncertified registered-family bootstrap outside the conformal certificate, and the
predictor probe is one model class on one feature family. Production traces, throughput benchmarks, and cross-operator exchangeability remain untested.

\section{Conclusion}
\label{sec:conclusion}

Safe search budgets are properties of a realized build, not only of a dataset and nominal recipe.
Across three dataset families, three implementations, and four scales, rebuilds cause large
finite-action variation and 17--24\% incremental transport risk. The theory separates
transcript-limited non-portability (premise measured) from fixed-target recovery (certification power
binds); the certified tiered pooling policy and the deterministic contract give the first measured
risk--cost frontier for the deployment decision. Verify determinism when feasible; otherwise pool with the conformal guarantee when history permits;
otherwise certify with declared margin, expecting only the maximum budget to clear; otherwise abstain.
The transport boundary is now measured at every layer we probe.

\section*{Ethics Statement}
This study uses only public benchmark datasets (SIFT, GloVe, deep-image, Arxiv-Nomic embeddings) and
synthetic workloads derived from them; it involves no human subjects, no personally identifiable data,
and no deployment on production systems. The reported risks concern internal quality-budget decisions in
ANN serving and do not involve protected attributes or fairness-sensitive outcomes. All evidence derives
from preregistered, forensically audited pipelines; no sealed evaluation roles were accessed without
preregistration.

\section*{Reproducibility Statement}
Every number in this paper traces to committed result tables, preregistered manifests (with commit
hashes and timestamps frozen before evaluation-role access), and deterministic test suites listed in
Appendix~\ref{app:repro}. The anonymous artifact contains the frozen result trees, analysis scripts,
build/search code, and checksums; the 1M and 10M cells carry their own preregistration manifests. All
conclusions are conditional on registered builds as documented in the claim registry
(Appendix~\ref{app:claims}).

\section*{AI Use Statement}
Generative AI tools assisted with organizing research artifacts, developing theory-experiment
crosswalks, checking document structure, drafting and editing manuscript text, and supporting code and
artifact validation. All mathematical statements, proofs, experimental results, citations, and
AI-assisted content require final author review, and the author assumes responsibility for the submitted
manuscript.

\begin{thebibliography}{99}
\itemsep0pt

\bibitem[Angelopoulos et~al.(2025)]{angelopoulos2025ltt}
Angelopoulos, A.~N., Bates, S., Candes, E.~J., Jordan, M.~I., and Lei, L.
Learn then test: Calibrating predictive algorithms to achieve risk control.
\emph{Annals of Applied Statistics}, 2025.

\bibitem[Angelopoulos et~al.(2024)]{angelopoulos2024crc}
Angelopoulos, A.~N., Bates, S., Fisch, A., Lei, L., and Schuster, T.
Conformal risk control.
\emph{International Conference on Learning Representations}, 2024.

\bibitem[Bae et~al.(2026)]{bae2026qbat}
Bae, J., Ham, T.~J., Li, A., Chockchowwat, S., and Papakonstantinou, Y.
QBAT: Model-based query budget autotuner for clustering-based approximate nearest neighbor search.
\emph{Proceedings of the VLDB Endowment}, 2026.

\bibitem[Bates et~al.(2021)]{bates2021rcps}
Bates, S., Angelopoulos, A.~N., Lei, L., Malik, J., and Jordan, M.~I.
Distribution-free, risk-controlling prediction sets.
\emph{Journal of the ACM}, 68(6), 2021.

\bibitem[Ben-David et~al.(2010)]{bendavid2010theory}
Ben-David, S., Blitzer, J., Crammer, K., Kulesza, A., Pereira, F., and Vaughan, J.~W.
A theory of learning from different domains.
\emph{Machine Learning}, 79:151--175, 2010.

\bibitem[Chatzakis et~al.(2025)]{chatzakis2025darth}
Chatzakis, M., Papakonstantinou, Y., and Palpanas, T.
DARTH: Declarative recall through early termination for approximate nearest neighbor search.
\emph{Proceedings of the ACM on Management of Data}, 3(4), 2025.

\bibitem[Duchi et~al.(2025)]{duchi2025multienv}
Duchi, J.~C., Gupta, S., Jiang, K., and Sur, P.
Predictive inference in multi-environment scenarios.
\emph{Statistical Science}, 40(3):392--416, 2025.

\bibitem[Elliott and Clark(2024)]{elliott2024order}
Elliott, O.~P. and Clark, J.
The impacts of data, ordering, and intrinsic dimensionality on recall in hierarchical navigable small
worlds. \emph{arXiv:2405.17813}, 2024.

\bibitem[El-Yaniv and Wiener(2010)]{elyaniv2010selective}
El-Yaniv, R. and Wiener, Y.
On the foundations of noise-free selective classification.
\emph{Journal of Machine Learning Research}, 11:1605--1641, 2010.

\bibitem[Fu et~al.(2019)]{fu2019nsg}
Fu, C., Xiang, C., Wang, C., and Cai, D.
Fast approximate nearest neighbor search with the navigating spreading-out graph.
\emph{Proceedings of the VLDB Endowment}, 12(5):461--474, 2019.

\bibitem[Garivier and Kaufmann(2016)]{garivier2016bai}
Garivier, A. and Kaufmann, E.
Optimal best arm identification with fixed confidence.
\emph{Proceedings of the 29th Conference on Learning Theory}, pages 998--1027, 2016.

\bibitem[Horchidan et~al.(2025)]{horchidan2025conann}
Horchidan, S., Zeiher, F., Bostrom, H., and Carbone, P.
ConANN: Conformal approximate nearest neighbor search.
\emph{Proceedings of the VLDB Endowment}, 19(1):29--42, 2025.

\bibitem[Johansson et~al.(2019)]{johansson2019support}
Johansson, F.~D., Sontag, D., and Ranganath, R.
Support and invertibility in domain-invariant representations.
\emph{Proceedings of the 22nd International Conference on Artificial Intelligence and Statistics}, pages
527--536, 2019.

\bibitem[Johnson et~al.(2021)]{johnson2021faiss}
Johnson, J., Douze, M., and Jegou, H.
Billion-scale similarity search with GPUs.
\emph{IEEE Transactions on Big Data}, 7(3):535--547, 2021.

\bibitem[Li et~al.(2020)]{li2020learned}
Li, C., Zhang, M., Andersen, D.~G., and He, Y.
Improving approximate nearest neighbor search through learned adaptive early termination.
\emph{Proceedings of the 2020 ACM SIGMOD International Conference on Management of Data}, pages
2539--2554, 2020.

\bibitem[Malkov and Yashunin(2020)]{malkov2020hnsw}
Malkov, Y.~A. and Yashunin, D.~A.
Efficient and robust approximate nearest neighbor search using hierarchical navigable small world graphs.
\emph{IEEE Transactions on Pattern Analysis and Machine Intelligence}, 42(4):824--836, 2020.

\bibitem[Prokhorenkova and Shekhovtsov(2020)]{prokhorenkova2020graph}
Prokhorenkova, L. and Shekhovtsov, A.
Graph-based nearest neighbor search: From practice to theory.
\emph{Proceedings of the 37th International Conference on Machine Learning}, pages 7803--7813, 2020.

\bibitem[Subramanya et~al.(2019)]{subramanya2019diskann}
Subramanya, S.~J., Devvrit, Simhadri, H.~V., Krishnaswamy, R., and Kadekodi, R.
DiskANN: Fast accurate billion-point nearest neighbor search on a single node.
\emph{Advances in Neural Information Processing Systems}, 32, 2019.

\bibitem[Wang et~al.(2026a)]{wang2026annie}
Wang, Z., Chatzakis, M., Wang, Q., Palpanas, T., Wang, P., and Wang, W.
ANNiE: A learned query cost estimator for graph-based approximate nearest neighbor search.
\emph{Proceedings of the VLDB Endowment}, 19(11):3820--3833, 2026.

\bibitem[Wang et~al.(2022)]{wang2022safebai}
Wang, Z., Wagenmaker, A., and Jamieson, K.
Best arm identification with safety constraints.
\emph{Proceedings of the 25th International Conference on Artificial Intelligence and Statistics}, 2022.

\bibitem[Xu et~al.(2022)]{xu2022maintenance}
Xu, Z., Zhao, W., Tan, S., Zhou, Z., and Li, P.
Proximity graph maintenance for fast online nearest-neighbor search.
\emph{arXiv:2206.10839}, 2022.

\bibitem[Yu(1997)]{yu1997assouad}
Yu, B.
Assouad, Fano, and Le Cam.
In \emph{Festschrift for Lucien Le Cam}, pages 423--435. Springer, 1997.

\bibitem[Zhang and Miller(2026)]{zhang2026exploration}
Zhang, C. and Miller, R.~J.
Distribution-aware exploration for adaptive HNSW search.
\emph{Proceedings of the ACM on Management of Data}, 4(1), 2026.

\end{thebibliography}


\appendix
\section{Related Work: Full Discussion}
\label{app:related}

\paragraph{Positioning.}
The main-text positioning paragraph cites the systems per line; the notes below add
differentiating context.
\label{sec:formulation}

Let $\Ecal$ be a preregistered finite set of build environments. An environment $E\in\Ecal$ consists of a
serialized index together with its implementation, construction configuration, input permutation, seed when
meaningful, toolchain, and replay semantics. Each implementation has its own finite native action set
$\Acal_E$, and queries follow a frozen law $q\sim P_Q$.

For action $a\in\Acal_E$, define the endpoint-aware failure indicator and implementation-internal cost by
\begin{equation}
Z_{E}(q,a)=\mathbf{1}\{\text{the returned result fails the registered quality target}\},\qquad
C_{E}(q,a)\ge 0,
\label{eq:zc}
\end{equation}
and the minimum observed safe budget by
\begin{equation}
B_{E}(q)=\min\{a\in\Acal_{E}:Z_{E}(q,a)=0\},\qquad
B_{E}(q)=\BOT\ \text{if no registered action succeeds}.
\label{eq:bmin}
\end{equation}
The symbol $\BOT$ is categorical and is never imputed as the largest action or a larger numerical budget. A
policy observes a transcript $T$ and query-side observables, then returns an action or abstains. Its
absolute target risk is
\begin{equation}
r_{E}^{\pi}=\Pr_{q\sim P_{Q}}\!\left[Z_{E}(q,\pi_{T,q})=1\right].
\label{eq:risk}
\end{equation}
We use risk tolerance $\delta=0.05$, certification error probability $\alpha=0.05$, and safety margin
$\gamma>0$ stated per experiment; the materiality gate $0.02$ is a preregistered scientific threshold---chosen, before any 100K evaluation
role was accessed, as the smallest risk level that would still be operationally material for budget
reuse---not a utility function. ``Preregistered'' in this paper means frozen in a versioned manifest with commit hash and
timestamp before the corresponding evaluation role was accessed; it does not refer to an external OSF
registration. A rejection event is $R$. An action $a_f$ is a fallback only if it has independent safety
evidence under the same target and failure event; otherwise rejection produces abstention.

The primary estimand of Section~\ref{sec:results} is the stage-registered endpoint-aware incremental
transport risk
\begin{equation}
\Delta_{s\rightarrow t}
=\Pr_{q\sim P_{Q}}\!\left[Z_{t}\big(q,B_{s}(q)\big){=}1\right]-r_{t}^{\mathrm{ref}},
\label{eq:delta}
\end{equation}
where, per sample, the transported action maps $B_s(q){=}\BOT$ to the maximum registered action (the
registered historical convention), and the reference event is the target-bottom failure
$Z_t(q,\underline{a}_t){=}1$ with $\underline{a}_t$ the smallest registered action on the target; for the
four HNSW cells this coincides with the reference event of the repaired $h{=}10$ common estimand. Intervals are 95\%
query-bootstrap over 5{,}000 resamples (seed 991) retaining each sampled query's full vector of directed-pair
observations, conditional on the registered build family. Heterogeneity is reported through two separated
variation families,
\begin{equation}
V_{\mathrm{fin}}=\Pr_{q}\!\left[\big|\{B_{E}(q):E\in\Ecal_{\mathrm{reg}},\,B_{E}(q)\neq\BOT\}\big|>1\right],
\qquad
V_{\mathrm{end}}=\Pr_{q}\!\left[\mathbf{1}\{B_{E}(q){=}\BOT\}\ \text{nonconstant over }E\right],
\label{eq:var}
\end{equation}
finite-action variation and endpoint-state variation; reported separately so response heterogeneity cannot
be attributed to unresolved endpoints. We distinguish registered-grid unresolved, true endpoint infeasible,
and right-censored queries. Current finite grids often cannot separate the latter two, so we report
unresolved mass without promoting it to true infeasibility.

\paragraph{Cost naming.}
ROM-NDC quantities in earlier drafts are renamed \emph{DistComp} throughout:
$\DistComp(x \text{ vs } y)$ is the ratio of mean numbers of distance computations of $x$ to $y$, computed
only within one implementation. ``NDC'' alone is avoided to prevent confusion with NDCG.



\section{Complete Notation and Semantic Contracts}
\label{app:notation}

\begin{table}[h]
\centering
\caption{Core notation and scope guardrails.}
\small
\begin{tabular}{lll}
\toprule
Symbol & Meaning & Scope guardrail \\
\midrule
$\Ecal$ & preregistered build-environment set & no population over unseen builds \\
$E$ & one serialized build environment & includes replay semantics \\
$\Acal_E$ & finite native action set & implementation specific \\
$Z_E(q,a)$ & endpoint-aware failure indicator & shared across roles \\
$B_E(q)$ & minimum observed safe action & registered grid only \\
$\BOT$ & unresolved registered endpoint & never numerically imputed \\
$C_E(q,a)$ & distance computations & within one implementation \\
$\Delta_{s\to t}$ & incremental transport risk (Eq.~\ref{eq:delta}) & registered-family conditional \\
$V_{\mathrm{fin}},V_{\mathrm{end}}$ & variation families (Eq.~\ref{eq:var}) & reported separately \\
$\delta,\alpha,\gamma,M$ & 0.05; 0.05; margin; frozen family size & preregistered \\
\bottomrule
\end{tabular}
\end{table}

For HNSW, raw fixed \texttt{efSearch} is treated as a finite action rather than an assumed monotone
stopping sequence. For Vamana-style search, $l$ and beam width are distinct; beam width is fixed to one in
the registered experiment. Native budget and cost counters remain implementation-specific.

\section{Proof of Theorem 1}
\label{app:proof1}

Define an induced binary decision $I$: return 0 when the chosen action lies in $D_0$, return 1 when it lies
in $D_1$, and assign actions outside both regions arbitrarily. Since $D_0\cap D_1=\varnothing$, an action
outside $D_i$ incurs loss at least $\Delta$. Therefore
\begin{equation}
\E_{0}[l_{0}]+\E_{1}[l_{1}]\ \ge\ \Delta\bigl\{P_{0}[I{=}1]+P_{1}[I{=}0]\bigr\}.
\label{eq:proof1}
\end{equation}
The minimum sum of binary testing errors equals $1-\TV$; the maximum expected loss is at least half the sum,
proving Theorem~1. Bretagnolle--Huber lower-bounds the error sum and Pinsker yields the KL consequence.
For adaptive observations the chain rule gives
\begin{equation}
K_{m}=\sum_{t=1}^{m}\E_{0}\!\left[\KL\bigl(P_{0}^{Y_{t}\mid H_{t-1},A_{t}}\,\bigl\|\,P_{1}^{Y_{t}\mid H_{t-1},A_{t}}\bigr)\right]\le m\,I^{\star},
\label{eq:chain}
\end{equation}
and data processing through the decision rule yields the stated testing consequence. The testing
inequalities are classical; only the Graph-ANNS deployment-loss reduction is domain specific. The
instantiation in Section~\ref{sec:scale} supplies the measured TV bound the premise requires.

\section{Proof of Theorem 2}
\label{app:proof2}

On the estimation event every retained action satisfies the screened bound. If a minimum-cost safe action
$a_k$ has margin $r_{a_k}\le\delta-2\varepsilon_R$, then
\begin{equation}
\hat r_{a_{k}}+\varepsilon_{R}\ \le\ r_{a_{k}}+2\varepsilon_{R}\ \le\ \delta,
\label{eq:retain}
\end{equation}
so it is retained; since the deployed action $\hat a$ minimizes estimated cost among retained actions,
\begin{equation}
\hat C_{\hat a}\ \le\ \hat C_{a_{k}}+\varepsilon_{C}\ \le\ C_{a_{k}}+2\varepsilon_{C}.
\label{eq:cost}
\end{equation}
If every other safe action costs strictly more than the optimum plus $2\varepsilon_C$, these inequalities
force exact selection on the event. Selection and certification roles are independent and the final action
is frozen; a valid upper certificate plus a union bound gives the $1-\alpha$ guarantee. On rejection the
procedure executes only an independently valid fallback; without one it abstains. The endpoint obstruction
(every in-family policy fails on the unresolved or infeasible subset), the identifiability obstruction
(Theorem~1), the no-margin obstruction (risks $\delta-\epsilon$ and $\delta+\epsilon$ become
indistinguishable as $\epsilon\downarrow0$), adaptive selection bias, and rejection without a safe fallback
are necessary obstructions, not an if-and-only-if theorem.

\section{Corollaries and Restricted Propositions}
\label{app:corollaries}

\paragraph{Reverse-pair symmetry.} If every unordered pair appears in both directions with equal weight,
each unequal pair contributes one under-budget and one over-budget direction, hence
\begin{equation}
P_{\mathrm{under}}=P_{\mathrm{over}}=\tfrac{1}{2}\,P\bigl[B_{s}\neq B_{t}\bigr],
\label{eq:sym}
\end{equation}
a combinatorial identity; asymmetric source-to-held-out-target designs are still needed for directional
portability claims.

\paragraph{Endpoint decomposition.} With $\eta_E$ the mass on which no registered action succeeds,
\begin{equation}
r_{E}^{\pi}=\eta_{E}+(1-\eta_{E})\,r_{E}^{\pi,\mathrm{feas}},
\label{eq:endpoint}
\end{equation}
so $\eta_E>\delta$ precludes in-family certification.

\paragraph{Failure complementarity.} With conditional rescue rate $\rho_S$,
\begin{equation}
r_{S}=r_{0}\,(1-\rho_{S}),
\label{eq:rescue}
\end{equation}
and marginal auxiliary accuracy cannot replace the conditional rescue rate; shared queues, visited sets,
truncation, or altered tie rules can invalidate the mechanical preservation condition.

\paragraph{Structural bridges.} The only generally valid bound retained is
\begin{equation}
\bigl|r_{G',a}-r_{G,a}\bigr|\ \le\ \Pr_{q}\{Z_{G}(q,a)\neq Z_{G'}(q,a)\},
\label{eq:bridge}
\end{equation}
and expansion-prefix arguments require fixed search semantics, preserved critical edges, priority margins,
and bounded frontier intruders; they do not imply a theorem for raw \texttt{efSearch} or $l$ without an
empirical implementation bridge.

\begin{figure}[t]
\centering
\includegraphics[width=.95\linewidth]{fig6_ablation.png}
\caption{Deterministic-contract ablation (descriptive chain). Left bars: minimum-safe-action variation per
tier; right line: median build time. Order pinning halves variation for free; single-threading removes all
variation and carries the build overhead.}
\label{fig:ablation}
\end{figure}

\begin{table}[t]
\centering
\caption{Six-cell economics matrix (selected rows; the full 36-cell matrix with per-cell NOT\_ESTIMABLE
reasons is in the artifact). N/E = not estimable from frozen evidence; never imputed.}
\label{tab:economics}
\footnotesize
\setlength{\tabcolsep}{3.4pt}
\begin{tabular}{lccccc}
\toprule
Component & hnS & hnA & FS & FA & Vam \\
\midrule
Build time (D3 median) & 12.9\,s & 46.6\,s & N/E & N/E & N/E \\
Profiling, resident 59q (8t) & 0.117\,s & 0.994\,s & 0.017\,s & 0.248\,s & N/E \\
Break-even, search-only & NO\_FINITE & 11{,}832\,q & N/E & N/E & N/E \\
Break-even, wall-clock & NO\_FINITE & 104{,}383\,q & N/E & N/E & N/E \\
Candidate replay / control & N/E & N/E & N/E & N/E & N/E \\
\bottomrule
\end{tabular}
\end{table}

\begin{table}[t]
\centering
\caption{Implementation-native semantics. Raw action values and distance-computation counts are never
compared across implementations.}
\label{tab:semantics}
\footnotesize
\setlength{\tabcolsep}{4pt}
\begin{tabular}{p{2.4cm}p{3.4cm}p{2.1cm}p{5.2cm}}
\toprule
Implementation & Registered build environment & Native action & Semantic guardrail \\
\midrule
hnswlib HNSW & seed $\times$ input order (24 builds/dataset) & \texttt{efSearch} & candidate-list control; not a DC cap \\
Faiss HNSW & registered input permutation (24 builds/dataset) & \texttt{efSearch} & implementation-specific; not equated \\
DiskANN3/Vamana-style & registered permutation; fixed options (12/dataset) & $l$, beam width $=1$ & list size; $\neq$ HNSW \texttt{efSearch} \\
\bottomrule
\end{tabular}
\end{table}

\begin{table}[t]
\centering
\caption{Theoretical results, assumptions, novelty boundary, and their empirical instantiation.}
\label{tab:theory}
\footnotesize
\setlength{\tabcolsep}{3.5pt}
\begin{tabular}{p{2.5cm}p{3.3cm}p{3.5cm}p{3.5cm}}
\toprule
Result & Core assumptions & Guarantee & Status of tools \\
\midrule
Transcript lower bound & conflicting regions; probe laws & four-outcome loss bound & classical testing; premise measured (\S\ref{sec:scale}) \\
Conditional recovery & feasibility, identifiability, & fixed-target familywise & new domain combination; margin \\
 & margin, evidence, fallback & safety + near-optimal cost & audited per target (\S\ref{sec:constructive}) \\
Certification scale & frozen $M$, margin $\gamma$ & sample order; zero-failure $m_{\min}$ & classical concentration \\
Economic gate & complete costs & finite break-even iff $D>0$ & accounting identity \\
\bottomrule
\end{tabular}
\end{table}

\begin{table}[t]
\centering
\caption{Contract outcomes and primitive profiling economics (8-thread resident measurements; primitive
operations only, not a deployment ledger).}
\label{tab:contract}
\footnotesize
\setlength{\tabcolsep}{3.5pt}
\begin{tabular}{p{3.7cm}llp{4.2cm}}
\toprule
Result & SIFT-100K & Arxiv-100K & Interpretation \\
\midrule
Deterministic transport risk & 0 & 0 & registered same-data/toolchain contract \\
Recall difference vs.\ original & $+0.00027$ & $-0.00040$ & within the registered $\pm0.001$ tolerance \\
Mean safe-budget ratio & 0.976 & 0.999 & descriptive common-feasible summary \\
p95 / p99 budget ratio & 1.00 / 1.00 & 1.00 / 1.00 & no measured tail change \\
Build-time ratio vs.\ original & $3.52\times$ & $1.71\times$ & determinism has offline overhead \\
59-query primitive profiling & H 0.117\,s; F 0.017\,s & H 0.994\,s; F 0.248\,s & cheap, operator dependent \\
\bottomrule
\end{tabular}
\end{table}

\begin{table}[t]
\centering
\caption{Incremental transport risk (Eq.~\ref{eq:delta} caliber) under relaxed quality thresholds, one
uniform estimator over all cells; the $h{=}10$ column reproduces the registered incremental estimates
(Appendix~F crosswalk).}
\label{tab:hsens}
\footnotesize
\setlength{\tabcolsep}{4pt}
\begin{tabular}{lccc}
\toprule
Cell & $h{=}8$ & $h{=}9$ & $h{=}10$ \\
\midrule
hnswlib SIFT-100K & 15.79\% & 19.51\% & 21.57\% \\
hnswlib Arxiv-100K & 9.35\% & 13.41\% & 17.12\% \\
Faiss SIFT-100K (clean) & 10.56\% & 16.77\% & 23.60\% \\
Faiss Arxiv-100K (clean) & 5.96\% & 10.31\% & 17.77\% \\
SIFT-1M (preregistered) & 19.72\% & 21.98\% & 21.87\% \\
\bottomrule
\end{tabular}
\end{table}

\begin{figure}[t]
\centering
\includegraphics[width=.9\linewidth]{fig3_crossfamily.png}
\caption{Cross-family evidence. Bars show safe-budget response variation; points show incremental transport
risk with 95\% query-bootstrap intervals. Vamana marks (daggered in Table~\ref{tab:vamana}) retain their
stage-specific estimand and are supporting rather than harmonized evidence.}
\label{fig:crossfamily}
\end{figure}

\begin{table}[t]
\centering
\caption{Preregistered P6 probes: scale replication, pooling boundary, predictor transfer, and transcript
distinguishability.}
\label{tab:p6}
\footnotesize
\setlength{\tabcolsep}{3pt}
\begin{tabular}{p{3.2cm}p{6.3cm}p{4.2cm}}
\toprule
Probe & Result & Boundary closed \\
\midrule
SIFT-1M cell (prereg.; 8 builds, 500 q) & incr.\ risk 21.87\% [21.52, 22.21]; $V_{\mathrm{fin}}$ 98.4\% vs $V_{\mathrm{end}}$ 13.2\%; contract byte-identical; profiling 0.067\,s & phenomenon replicates at $10\times$ scale (100K: 21.55\%) \\
Pooling at 1M & $k{=}7$ leaves 10.4\% risk (100K w/ 22 sources: $<1\%$); quantile pooling: $q{=}0.9 \approx$ max, median unusable & pooling requirement scale-dependent \\
Learned predictor (HistGB, ef{=}10 transcript) & transfer 23.9\%/21.2\% vs naive 21.1\%/17.3\%; in-target CV 24.1\%/22.2\%; $R^2=0.015$ & layer-3 non-portability; $B(q)$ build-structure dominated \\
Distinguishability (552 pairs, $k\in\{1,3,6\}$) & hits: acc.\ 0.49--0.51, TV $\le0.25$; runtime features: acc.\ 0.56--0.58 (max 0.72) & Thm.~1 premise holds for hit transcripts; runtime probes partially separate \\
\bottomrule
\end{tabular}
\end{table}

\begin{figure}[t]
\centering
\includegraphics[width=.95\linewidth]{fig4_pooling.png}
\caption{Robust source pooling: transport risk versus number of pooled source builds (mean and p95 over 50
draws; log scale). $k{=}1$ reproduces the registered naive-transport risks; $k\ge10$ falls below 2.4\% at
100K. At SIFT-1M the decay is slower ($k{=}7$ still leaves 10.4\%): the pooling requirement grows with
scale.}
\label{fig:pooling}
\end{figure}

\begin{figure}[t]
\centering
\includegraphics[width=.75\linewidth]{fig5_plane.png}
\caption{Constructive decision plane (hnswlib). Risk versus DistComp against the per-query oracle;
horizontal gates at 2\% (materiality) and 5\% (risk tolerance). M2 occupies the low-risk/low-cost corner
without certification; certified M3 coincides with M5.}
\label{fig:plane}
\end{figure}

\begin{table}[t]
\centering
\caption{End-to-end ledger per target build (hnswlib; offline cost measured, risk as in
Table~\ref{tab:constructive}). ``Certified'' = zero-failure CP at the stated $m$.}
\label{tab:e2e}
\footnotesize
\setlength{\tabcolsep}{3pt}
\begin{tabular}{p{3.1cm}p{2.6cm}p{2.6cm}p{1.4cm}p{2.4cm}}
\toprule
Route & Offline cost / target & Risk (SIFT/Arxiv) & DistC. & Certified \\
\midrule
M1 naive transport & 0 & 21.9\%/18.1\% & 2.6/2.7$\times$ & no \\
M2 pooled replay & 0 (needs label cache) & 0.18\%/0.09\% & 1.45/1.34$\times$ & no \\
Profile $+$ certify (59\,q, truth incl.) & 0.117\,s / 0.994\,s & 0.87\%/1.36\% & 4.16/5.28$\times$ & max-budget only; $P{=}0.47$--$0.63$ \\
Profile $+$ select $+$ certify ($m{=}94$) & $\le$0.8\,s / $\le$3\,s & --- & --- & 0/48 pass \\
Always max (M5) & 0 & 0.80\%/1.26\% & 4.09/5.14$\times$ & no \\
Deterministic contract (D3) & 12.9\,s / 46.6\,s build & incr.\ 0 (abs.\ $=$ source) & 0.98/0.999$\times$ & by replay identity \\
\bottomrule
\end{tabular}
\end{table}

\section{Complete Experimental Protocol}
\label{app:protocol}

The registered unit is a serialized build, not a query row. The strict HNSW evidence uses 24 hnswlib and 24
clean Faiss-100K builds per dataset and all 552 non-diagonal directed pairs; the Vamana-style analysis uses
12 builds and 36 preregistered directions; the 1M cell uses 8 builds, 56 directed pairs, and 500
preregistered evaluation queries. Every 100K cell contains 750 evaluation queries at Recall@10 with $h=10$.
The clean Faiss query roles have zero ID, raw-content, and normalized-content overlap with the indexed base
and historical roles; the Vamana stages carry the same audit at zero overlap with the caveats stated in
Section~\ref{sec:protocol}. Native action grids and counters are never numerically equated across
implementations.

The query bootstrap resamples query identities while retaining the complete vector of registered
directed-pair observations; it quantifies query-distribution uncertainty conditional on the finite build
family. Leave-one-build-out and leave-one-seed/order analyses test domination but do not define an
unseen-build population. Top-contributor deletion removes 1\% of evaluation queries under a frozen
contribution rule (eight queries under the registered hnswlib ordering). The confidence intervals support
registered-family conclusions only. An unresolved endpoint remains $\BOT$ in all categorical analyses.
Mean, p95, p99, build cost, truth cost, control cost, and fallback cost are separate fields; missing items
remain not estimable.

\begin{table}[t]
\centering
\caption{Vamana-style boundary panel (dagger: stage-specific estimand, beam width one, asymmetric
directions; not harmonized with the HNSW estimand and not merged into its range).}
\label{tab:vamana}
\footnotesize
\setlength{\tabcolsep}{2.6pt}
\begin{tabular}{llccccc}
\toprule
Impl. & Dataset & B/p & Cat.\ var. & Risk [95\% CI] & Cens. & Cost (descr.) \\
\midrule
DiskANN3/Vam.$^\dagger$ & SIFT-100K & 12/36 & 71.07\% & 16.86\% {\tiny[15.65, 18.10]} & 0.27\% & $\sim$1.00\% {\tiny[$-$0.94, 3.04]} \\
DiskANN3/Vam.$^\dagger$ & Arxiv-100K & 12/36 & 48.27\% & 11.37\% {\tiny[10.20, 12.60]} & 0.80\% & $\sim$1.11\% {\tiny[$-$1.05, 3.22]} \\
\bottomrule
\end{tabular}
\end{table}

\begin{figure}[t]
\centering
\includegraphics[width=.8\linewidth]{fig2_workflow.png}
\caption{ICBA workflow. The procedure validates semantics before estimating transport, separates selection
from certification, and treats fallback and abstention as distinct terminal outcomes.}
\label{fig:workflow}
\end{figure}

\section{Cross-Family Details and Estimand Crosswalk}
\label{app:crossfamily}

In hnswlib, same-seed cross-order inconsistency is 57.43\% (SIFT) and 45.23\% (Arxiv), substantially larger
than same-order cross-seed inconsistency (13.25\% and 12.62\%); the comparison is descriptive because
factors are not independently randomized across implementations. Faiss does not expose a reliable
construction seed under the registered interface, so fixed input permutations define its 24 build
environments; the clean 100K rerun removes query/base self-matches. DiskANN3/Vamana-style preflight
establishes fixed-configuration replay and nonzero differences across registered permutations, but the
frozen Vamana event table cannot reconstruct the harmonized HNSW source-action estimand; Vamana results
remain stage specific. Finite-action and endpoint variation are reported separately (75.87--91.07\% versus
0.53--2.67\%); unresolved mass is 0.072/0.167\% (Faiss SIFT/Arxiv) and 0.767/1.067\% (hnswlib SIFT/Arxiv).

\paragraph{Estimand crosswalk (hnswlib).}
The E4 primary oracle-transport estimand gives 21.57\% [20.74, 22.40] (SIFT) and 17.12\% [16.20, 18.06]
(Arxiv) under crossed seed-query bootstrap; the original cross-family stage gives 21.57\% [20.77, 22.35]
and 17.12\% [16.24, 18.00]; the repaired $h{=}10$ hit-count common estimand---adopted in
Table~\ref{tab:crossfamily} because the clean Faiss evidence exists only under hit-count semantics---gives
21.55\% [20.76, 22.33] and 17.17\% [16.33, 18.06]. All three coexist as registered values; the paper adopts
the repaired common estimand for the four HNSW cells and does not modify any frozen number.

\section{Falsification-Ladder Casebook}
\label{app:ladder}

\begin{table}[t]
\centering
\caption{Four recovery families and the first prerequisite falsified by sealed evidence; the 16-route
casebook is in Appendix~\ref{app:ladder}.}
\label{tab:ladder}
\footnotesize
\setlength{\tabcolsep}{3pt}
\begin{tabular}{p{2.3cm}p{4.2cm}p{4.2cm}p{2.7cm}}
\toprule
Recovery family & Strongest positive evidence & Decisive limitation & Inference \\
\midrule
Reuse / recalibration & fixed-target safety in selected designs; & no stable history increment; & fixed-target safety only \\
 & probe quantifies the negative side ($R^2{=}0.015$) & disjoint eval.\ removes gain & \\
Build selection / stable constr. & safe-selection theory valid; isolated gains & no CI-supported joint action (72 frozen) & generator closed \\
Multi-lane / portal rescue & complementarity replicates after ID repair & no truth-free trigger; 79.5--93.4\% compression & mechanism only \\
Certified auditor / fallback & 36/36 accepted decisions safe & collapse to max budget; regret 470.3 / 684.5 & diagnostic, not optimizer \\
\bottomrule
\end{tabular}
\end{table}

\begin{figure}[t]
\centering
\includegraphics[width=.95\linewidth]{fig7_ladder.png}
\caption{Theory-guided falsification ladder. Each recovery family is pruned at its first unsupported
prerequisite. Partial mechanism evidence is retained without being promoted to a deployable algorithm.}
\label{fig:ladder}
\end{figure}

Target-only residual recalibration initially reduced mean distance computations by 29.40\% (SIFT) and 26.17\%
(Arxiv) in a paired fixed-target design; history-assisted variants added only 0.90\% [$-8.73$, 11.25] and
$-0.13$\% [$-9.24$, 10.26]; under strict disjoint-query cross-fitting the target-only effects became
$-7.35$\% and $-12.49$\%. Build selection failed earlier: among 72 frozen build-budget actions, SIFT had no
point-estimate-or-interval supported jointly feasible action and Arxiv had one point-estimate action with no
interval support. One protected-edge repair preserved graph invariants yet reduced Recall@10 by 0.002 and
worsened both endpoint families. Fixed portal sets showed genuine rescue complementarity (1{,}064 and 727
additional hits; 426 and 348 threshold rescues) but cost $2.5$--$2.8\times$ a matched native lane and
remained 3.3--4.8 recall points lower; eight preregistered truth-free triggers showed no positive validation
advantage, and a shared-frontier bound implies 79.5--93.4\% auxiliary-cost compression before viability
(provenance for these aggregates is itemized in the artifact; three of them are recorded as
origin-unlocated and are being re-derived from row-level records). Finally, the conservative auditor
demonstrated safety without value across 60 retrospective decisions (4 direct reuses, 32 certified
fallbacks, 24 abstentions; all 36 accepted decisions safe, every accepted action at the maximum registered
budget).

\begin{table}[h]
\centering
\caption{Sixteen route-level tests summarized by family and first failed condition. Three historical
aggregates (the recalibration percentages, the protected-edge repair cost pair, and the exact matched-lane
ratios) lack stored provenance and are flagged for re-derivation in the artifact.}
\small
\begin{tabular}{lll}
\toprule
Family & Route & First failed prerequisite \\
\midrule
Reuse/recal. & source-only global reuse & nontrivial portability \\
Reuse/recal. & target-only residual calibration & disjoint-query replication \\
Reuse/recal. & history-assisted calibration & stable history increment \\
Reuse/recal. & adaptive sentinel allocation & matched-risk incremental value \\
Selection & certified build selection & jointly feasible candidate \\
Selection & portfolio racing & Stage-I feasibility \\
Selection & protected-edge repair & native recall/budget response \\
Selection & behavioral operator search & Pareto improvement \\
Selection & critical-frontier stabilization & structure-to-budget bridge \\
Portal & fixed portal union & absolute safety and cost \\
Portal & per-query portal oracle & deployable observability \\
Portal & truth-free trigger & positive validation advantage \\
Portal & shared frontier & realizable cost compression \\
Auditor & ordered certification & native-action nesting semantics \\
Auditor & certified fallback ladder & useful intermediate fallback \\
Auditor & unified ICBA decision & decision regret and economics \\
\bottomrule
\end{tabular}
\end{table}

The ladder preserves three gaps: an oracle action may exist without an observable signal; an observable
predictor may lack certification margin; and a certifiable action may have no mean, tail, or break-even
value. Candidate attainability precedes all three.

\section{Claim Registry}
\label{app:claims}

The final theory removes several overclaims: raw \texttt{efSearch} is not assumed to generate nested
population failure events; structural similarity does not guarantee budget stability; query bootstrap is
not build-level inference; a maximum registered action is not a safe fallback without evidence; mean
savings do not imply tail savings; the negative result is transcript limited rather than an absolute
hidden-build theorem; neither main theorem is a new general Le Cam, KL, Learn-Then-Test, or best-arm
result.

\paragraph{Permitted claims.}
Registered-family HNSW evidence at 100K and the preregistered 1M replication; stage-specific supporting
Vamana evidence; fixed-target certification semantics; the deterministic same-environment contract with its
chain ablation; robust source pooling as an uncertified registered-family baseline with measured scale
dependence; grid-robustness at phenomenon level; transcript-level indistinguishability for the registered
probe class; operator-dependent cost; exchangeable-build conformal pooling with finite-sample per-query
validity (Theorem~3), valid at every non-vacuous $(\alpha,k)$ on both implementations, under the stated
build-exchangeability assumption; tiered composition risk bound at measured repeat rates.

\paragraph{Prohibited claims.}
Universal Graph-ANNS behavior; open-world recovery; unseen-build safety (build exchangeability is an
assumption, not a certification, of the deployment process); native-budget numerical equivalence;
universal profiling or conservative-cost conclusions; production/SOTA performance of the pooling policy
(it is a budget-selection layer, not a retrieval algorithm, and is not compared against SOTA indexes);
resolution-invariant risk levels; learned-predictor transfer beyond the probed feature class; conformal
validity outside the probed transcript/feature classes or for target-dependent source selection.

\section{Reproducibility and Artifact Manifest}
\label{app:repro}

The final evidence derives from frozen hnswlib, clean Faiss HNSW, and DiskANN3/Vamana-style result trees,
plus the preregistered SIFT-1M cell whose manifest (design, seeds, orders, grid, roles, stop rules) was
committed before any build or search; the forensics gate passed before any risk computation. Each stage
records implementation version, build configuration, input ordering, query-role hashes, native action
grid, estimand, confidence procedure, robustness deletions, and output checksums. The clean Faiss rerun
verifies zero query/base content overlap and 48/48 frozen index hashes; the deterministic contract verifies
byte-identical indexes and search outputs at 100K and 1M. The historical 123-entry baseline remains
conditionally reproducible at 111 matching entries, eight content differences, and four unavailable
unversioned cache files. The anonymous artifact omits identifying repository links and local paths; all
revision-cycle numbers trace to committed result tables with deterministic test suites (16, 24, 27, 16, 7,
22, and 17 checks per phase). Remaining code-audit boundaries---primitive rather than complete profiling
cost, hash-only role-access status, and the three origin-unlocated historical aggregates---are recorded
explicitly and are not used to strengthen the scientific claims.

\paragraph{Reproducibility statement.}
The manuscript reports registered build families, native action grids, query-role separation, endpoint
semantics, uncertainty units, and deterministic replay checks. The anonymous artifact contains scripts,
compact result tables, manifests, and checksums without author-identifying paths. No future-replication
truth or outcome contributes to the reported analyses; any metadata or hash-only access is separately
disclosed in the artifact audit.

\paragraph{AI use statement.}
Generative AI tools assisted with organizing research artifacts, developing theory--experiment crosswalks,
checking document structure, drafting and editing manuscript text, and supporting code and artifact
validation. All mathematical statements, proofs, experimental results, citations, and AI-assisted content
require final author review, and the author assumes responsibility for the submitted manuscript.

\end{document}



